% 第4章：零训练错误阈值区间随样本增多被夹紧。
% 教学分布：x ~ Uniform[0,1]，y = 1{x >= 0.6}。三行是构造的示意样本（非随机抽样结果）：
%   n=4 : 正常 0.12, 0.45；垃圾 0.71, 0.88            -> 区间 (0.45, 0.71]
%   n=8 : 正常 0.08, 0.27, 0.43, 0.55；垃圾 0.66, 0.74, 0.83, 0.95 -> 区间 (0.55, 0.66]
%   n=16: 正常 0.04,0.11,0.19,0.26,0.33,0.41,0.49,0.58；垃圾 0.62,0.67,0.72,0.79,0.84,0.90,0.93,0.98 -> (0.58, 0.62]
% 竖带为 [0.6-eps, 0.6+eps]，eps = 0.05；风险 |t-0.6| > eps 恰为阈值落在带外。
% 坐标：分数 s 映射到 x = 14 + 88 s（mm）。
\begingroup
\definecolor{FigureBlue}{HTML}{7188A3}\definecolor{FigureBlueFill}{HTML}{E5EAF0}
\definecolor{FigureGreen}{HTML}{829C8B}\definecolor{FigureGreenFill}{HTML}{E8EEE9}
\definecolor{FigureYellow}{HTML}{B4A16F}\definecolor{FigureYellowFill}{HTML}{F1ECDD}
\definecolor{FigureRed}{HTML}{A7787D}\definecolor{FigureRedFill}{HTML}{F0E5E6}
\definecolor{FigureInk}{HTML}{3B4252}\definecolor{FigureMuted}{HTML}{6B7280}
\begin{tikzpicture}[figure picture,x=1mm,y=1mm]
\def\s#1{\fpeval{14+88*#1}}
% epsilon band around the true threshold
\fill[FigureGreenFill] (\s{0.55},-1) rectangle (\s{0.65},33);
\draw[draw=FigureGreen,line width=.8pt] (\s{0.6},-1)--(\s{0.6},33);
\node[figure note,anchor=south] at (\s{0.6},33.3) {真实分界 $0.6$ 及其两侧 $\varepsilon=0.05$};
% one row: #1 = y, #2 = label, #3 = negatives, #4 = positives, #5 = left end, #6 = right end
\newcommand{\squeezerow}[6]{%
  \draw[draw=FigureStroke,line width=.6pt] (\s{0},#1)--(\s{1},#1);
  \fill[FigureYellow] (\s{#5},{#1-1.1}) rectangle (\s{#6},{#1+1.1});
  \foreach \v in {#3}{\draw[draw=FigureBlue,fill=FigureBlueFill,line width=.7pt] (\s{\v},{#1+3}) circle(1.1);}
  \foreach \v in {#4}{\draw[draw=FigureRed,fill=FigureRed,line width=.7pt] (\s{\v},{#1+3}) circle(1.1);}
  \node[anchor=east] at (\s{0}-1.5,#1+1.5) {#2};
  \node[anchor=west,font=\figurefont\scriptsize] at (\s{1}+1.5,#1+1.5) {$(#5,#6]$};}
\squeezerow{25}{$n=4$}{0.12,0.45}{0.71,0.88}{0.45}{0.71}
\squeezerow{14}{$n=8$}{0.08,0.27,0.43,0.55}{0.66,0.74,0.83,0.95}{0.55}{0.66}
\squeezerow{3}{$n=16$}{0.04,0.11,0.19,0.26,0.33,0.41,0.49,0.58}{0.62,0.67,0.72,0.79,0.84,0.90,0.93,0.98}{0.58}{0.62}
\foreach \v in {0,0.2,0.4,0.6,0.8,1}{\node[anchor=north,font=\figurefont\scriptsize,text=FigureMuted] at (\s{\v},-1.5) {$\v$};}
\node[anchor=west,font=\figurefont\scriptsize,text=FigureMuted] at (\s{1}+1.5,-3.2) {零错误区间};
\end{tikzpicture}
\endgroup
